\section{Dénombrabilité}
\subsection{Définitions}

\begin{definition}{Dénombrable}{}
    Soit $I$ un ensemble. On dit que $I$ est \emph{dénombrable} 
    s'il existe une bijection entre $\N$ et $I$.
    On dira que $I$ est \emph{équipotent} à $\N$ (LHP).
\end{definition}

\begin{exemple}
    \begin{itemize}
        \item L'ensemble $\N$ est dénombrable (la bijection est l'identité).
        \item L'ensemble $\N^*$ est dénombrable (la bijection est $f : n \mapsto n+1$).
        \item $2 \N$ est dénombrable (la bijection est $f : n \mapsto 2n$).
        \item $\Z$ est dénombrable (la bijection est $f : n \mapsto p \text{ ou } -p-1$ selon la parité de $n$).
    \end{itemize}
\end{exemple}


\begin{definition}{Au plus dénombrable}{}
    Soit $I$ un ensemble. On dit que $I$ est \emph{au plus dénombrable} si $I$ est 
    fini ou dénombrable.
\end{definition}

\subsection{Cas des parties de \texorpdfstring{$\N$}{N}}

\begin{theoreme}{Nature des parties infinies de $\N$}{}
    Toute partie infinie de $\N$ est dénombrable.
\end{theoreme}

\idee{
    On numérote les éléments de $A$ par ordre croissant en construisant une 
    bijection et des récurrences.
}

\begin{proof}
    Soit $A \subset \N$ infinie. 

    Construisons par récurrence $\varphi : \N \to A$ en posant : $\varphi(0) = \min(A)$ car $A \neq \emptyset$.

    Et par récurrence, $\forall n \in \N, \, \varphi(n+1) = \min(\underbrace{A \setminus \underbrace{\{\varphi(0), \ldots, \varphi(n)\}}_{\text{est fini}}}_{\text{infini donc non vide}})$. \\

    \uline{Montrons que $\varphi$ est bijective :}

    \svspace On a $\varphi(n) = \min(\underbrace{A \setminus \{\varphi(0), \ldots, \varphi(n-1)\}}_{= B })$ 
    et $\varphi(n + 1) = \min(B \setminus \{ \varphi(n) \})$ 

    Donc $\forall k \in B, \, \varphi(n) \leqslant k$ et $\exists k_1 \in B \setminus \{\varphi(n)\}, \, \varphi(n + 1) = k_1$

    Donc $\varphi(n) \leqslant \varphi(n + 1)$ et $\varphi (n) \neq \varphi(n + 1)$ par construction.

    Donc $\varphi(n) < \varphi(n + 1)$, donc $\varphi$ est injective. \\

    Soit $a \in A$, alors $\{k \in A | k < a\}$ est fini et $(\varphi(n))_{n \in \N}$ est strictement croissante, donc
    $\exists n_0 \in \N, \, \varphi(n_0) \geqslant a$.

    Soit $n_0$ le plus petit entier tel que $\varphi(n_0) \geqslant a$.

    On a donc $\varphi(n_0 - 1) < a \leqslant \varphi(n_0)$.
  
    Donc $a = \varphi(n_0)$ par construction, donc $\varphi$ est surjective.
\end{proof}

\begin{proposition}{Nature des parties de $\N$}{}
    Dans $\N$, il y a deux types de parties : 
    \begin{itemize}
        \item Les parties majorées qui sont finies.
        \item Les parties non majorées qui sont infinies et dénombrables.
    \end{itemize}

    Toute partie de $\N$ est donc au plus dénombrable. 
\end{proposition}


\begin{proposition}{Caractérisation des ensembles au plus dénombrables}{}
    Soit $E$ un ensemble non vide. Alors $E$ est au plus dénombrable ssi il est en 
    bijection avec une partie de $\N$.
\end{proposition}


\begin{proof}
    $\Leftarrow$ Si $E$ est en bijection avec $I \subset \N$, alors si $I$ est finie, 
    $E$ l'est également, et si $I$ est infinie, alors $I$ est dénombrable donc $E$ aussi.

    $\Rightarrow$ Si $E$ est au plus dénombrable, alors si $E$ est fini, il est en bijection 
    avec $\ninter 1 {\text{Card}(E)} \subset \N$.
    S'il est dénombrable, par définition il est en bijection avec $\N$.
\end{proof}


\subsection{Opérations}

\begin{theoreme}{}{}
    $\N^2$ est dénombrable.
\end{theoreme}

\begin{proof} (non exigible)
    Considérons $\fonction \varphi {\N^2} \N {(p, q)}{2^p (2q + 1) - 1}$.

    $\varphi$ est bien définie, et si $a \in \N$, alors $a + 1 \in \N^*$, donc possède 
    une unique décomposition en facteurs premiers qui fournit $p$ et $q$ en isolant 
    la contribution de $2$ et le produit impair des autres premiers. 
\end{proof}

\begin{proposition}{Opérations conservant le caractère dénombrable}{}
    \begin{itemize}
        \item $\forall p \in \N, \, \N^p$ est dénombrable.
        \item Tout produit cartésien fini d'ensembles dénombrables est dénombrable.
    \end{itemize}

\end{proposition}

\idee{Par récurrence.}

\begin{proof}
    Par récurrence : 
    \begin{itemize}
        \item $\N^1$ est dénombrable 
        \item Supposons que $\N^p$ est dénombrable pour un certain $p$. Cela fournit $\varphi : \N \to \N^p$ bijective. 
    \end{itemize}

    Considérons $\psi : \N^2 \to \N$ bijective. Posons alors $\fonction f {\N^{p + 1}} \N {(n_1, \cdot, n_{p + 1})} {\psi(\varphi^{-1}(n_1, \ldots, n_p), n_{p + 1})}$.

    Alors $f$ est bijective. 

    \lvspace Soient $E_1, \ldots, E_n$ des ensembles dénombrables.

    Posons pour tout $i, \varphi_i : E_i \to \N$ bijective et $f : \N^p \to \N$ bijective.

    Alors $\fonction{}{E_1 \times \cdot \times E_p} \N {(e_1, \ldots, e_p)} {f(\varphi_1(e_1), \ldots, \varphi_p(e_p))}$ est bijective.
\end{proof}

\begin{theoreme}{}{}
    $\Z$ est dénombrable.
\end{theoreme}

\begin{proof}
Posons $\varphi : x \mapsto \frac n 2$ si $n \in 2 \N$ et $-(\frac {n+1} 2)$ sinon.

Alors $\varphi$ est injective et surjective dans $\Z$ donc elle est bijective.
Donc $\Z$ est équipotent à $\N$ donc $\Z$ est dénombrable.
\end{proof}

\begin{theoreme}{}{}
    $\Q$ est dénombrable. 
\end{theoreme}

\idee{
    Bijection qui associe à chaque rationnel le couple (numérateur, dénominateur)
    irréductible.
}

\begin{proof}
    Considérons $\fonction \varphi \Q {\Z \times \N^*}{\underbrace{ x = \frac p q}_{\text{irréductible}} } {(p, q)}$.

    $\varphi$ est bien définie et injective, car si $\varphi(x) = (p, q)$, alors $x = \frac p q$

    Mais $\Z \times \N^*$ est dénombrable, donc $\varphi(\Q) \subset \Z \times \N^*$ est 
    au plus dénombrable. 

    Or $\varphi(\Q)$ est en bijection avec $\Q$, donc $\Q$ est au plus dénombrable, or $\Q$ est 
    infini, donc $\Q$ est dénombrable.
\end{proof}

\begin{theorement}{Nature d'une union au plus dénombrable d'ensemble au plus dénombrables}{}
    Toute union \textbf{finie ou dénombrable} d'ensembles au plus dénombrables est au plus dénombrable. 
\end{theorement}

\begin{proof} (non exigible)\\
    Soit $I$ non vide au plus dénombrable.
    
    Pour $i \in I, \, I \subset \N$, il existe $P_i \subset \N$ non vide et $\fonction {\varphi_i} {P_i} {E_i} n {\varphi_i(n) }$ bijective. 

    Considérons $\fonction {\psi_i} {E_i} \N x {\varphi_i^{-1}(x)}$ alors $\psi_i$ n'est pas forcément 
    surjective mais elle est injective. 

    Considérons maitenant $\fonction {\psi} {\bigcup_{i \in I} E_i} {\N^2} x {(i, \psi_i(x))}$ où $i = \min \{ \alpha \in \N, \, x \in E_\alpha \}$.

    $\psi$ est bien définie, et $\psi$ est injective, car si $x = y = (i, j)$, 
    on a $x, y \in E_i$ et $j = \psi_i(x) = \psi_i(y)$, donc $x = y = \varphi_i(j)$.

    Donc $\displaystyle \bigcup_{i \in I} E_i$ est en bijection avec son image par $\psi$
    qui est une partie de $\N^2$, qui est donc au plus dénombrable.

\end{proof}

\subsection{Ensembles non dénombrables}

\begin{theoreme}{Nature de $\R$}{}
    $\R$ est non dénombrable.
\end{theoreme}

\begin{proof} (non exigible)
    Par l'absurde, supposons que $\varphi : \N \to \R$ est bijective. 

    Comme $\fonction{} \R {]0, 1[} x {\frac 1 {1 + e^{- x}}}$ est bijective, donc 
    il existe une bijection de $\N$ dans $]0, 1[$.

    Considérons pour $i \in \N^*$ la $i + 1^{\text{eme}}$ décimale du développement décimal 
    propre de $\varphi(i)$. Notons la $a_i$. 

    Posons $b_i = 4$ si $a_i \neq 4$ et $b_i = 5$ sinon, et posons $\displaystyle y = \sum_{i = 1}^\infty \frac{b_{i + 1}}{10^i} \in ]0, 1[$ 

    Posons $j = \psi^{-1}(y)$,
    la jème décimale de $y$ est $b_j$ par définition de $y$ sa $j$e décimale est 
    $a_j$, absurde. Donc $\psi \circ \varphi$ non bijective.
\end{proof}

\begin{theoreme}{Non équipotence de $E$ et $\mathscr P(E)$ - LHP}{}
    Soit $E$ un ensemble non vide. Alors il n'existe pas de bijection entre $E$ et $\mathscr P(E)$.
\end{theoreme}

\begin{proof}
    Par l'absurde, si $\fonction \psi E {\mathscr P(E)} x {\psi(x)}$ est bijective,

    Considérons $B = \{ x \in E | x \notin \psi(x) \} $ et $x_0 = \psi^{-1}(B)$.

    Si $x_0 \in B, \, x_0 \notin \psi(x_0) = B$ absurde.

    Si $x_0 \notin B, \, x_0 \in \psi(x_0)$ donc $x_0 \in B$ absurde. 
\end{proof}

\begin{remarque}
    On montrerait que $\mathscr P(\N)$ est en bijection avec $\R$.
\end{remarque}

\section{Familles sommables}
\subsection{Définition et premières propriétés}

\begin{definition}{Famille sommable}{}
    Soit $I$ un ensemble non vide et $(u_j)_{j \in I}$ une famille de complexes indexée par $I$. 
    On dit que cette famille est sommable ssi l'ensemble 
    $$\left\{\sum_{j \in J} |u_j|, \, J \subset I \text{ fini} \right\}$$
    est majoré. 
\end{definition}

\begin{proposition}{Stabilité par combinaison linéaire}{}
    Soit $(u_j)_{j \in I}$ une famille sommable de complexes et $I' \subset I$ non vide.
    \begin{itemize}
        \item Alors $(u_j)_{j \in I'}$ est sommable. 
        \item Si $(v_j)$ est une autre famille sommable, et $\alpha, \beta \in \C$, alors 
        $(\alpha u_j + \beta v_j)_{j \in I}$ est sommable.
    \end{itemize}
\end{proposition}

\idee{Majorer les sommes partielles.}

\begin{proof}
    Par hypothèse, $\exists M_1, M_2 > 0, \, \forall J \subset I$ fini, 

    $\displaystyle \sum_{j \in J} |u_j| \leqslant M_1$ et $\displaystyle \sum_{j \in J} |v_j| \leqslant M_2$.

    Soit $J \subset I'$ fini, alors $J \subset I$ donc $\sum_{j \in J } |u_j| \leqslant M_1$.

    Donc $(u_j)_{j \in I'}$ est sommable.

    Soit $J \subset I$ fini, alors
    $$\sum_{j \in J} |\alpha u_j + \beta v_j | \leqslant \sum_{j \in J} (|\alpha| |u_j| + |\beta| |v_j|) \leqslant |\alpha| M_1 + |\beta| M_2$$
\end{proof}

\subsection{Somme d'une famille sommable}

\begin{proposition}{Inégalité de comparaison}{}
    Soit $(u_j)_{j \in I}$ une famille sommable de complexes et $(v_j)_{j \in I}$
    une autre famille telle que $\forall j \in I, \, |v_j| \leqslant |u_j|$.

    Alors $(v_j)_{j \in I}$ est sommable.
\end{proposition}

\begin{proof}
    Soit $J \subset I$ fini, alors :
    $$\sum_{j \in J} |v_j| \leqslant \sum_{j \in J} |u_j| \leqslant M$$
\end{proof}


\begin{definition}{Somme d'une famille sommable}{}\\
    Soit $(u_i)_{i \in I}$ une famille sommable de réels positifs. 

    Alors $\displaystyle \left\{ \sum_{i \in J} u_i, \, J \subset I \text{ fini}  \right\}$ est une partie 
    de $\R$. 

    Sa borne supérieure est appellée \emph{somme} de la famille et notée $\displaystyle \sum_{i \in I} u_i$.
    Ainsi : $$\sum_{i \in I} u_i = \sup \left\{ \sum_{i \in J} u_i, \, J \subset I \text{ fini}  \right\}$$
\end{definition}

\begin{theoreme}{Somme d'une famille sommable de réels quelconques}{}
    Soit $(u_j)_{i \in I}$ une famille de réels. 
    Posons $\forall i \in I : $ 
    \begin{align*}
        u_i^+ = \frac 1 2 (|u_i| + u_i)\\
        u_i^- = \frac 1 2 (|u_i| - u_i)
    \end{align*}

    Alors $(u_j)_{j\in I}$ est sommable ssi $(u_j^+)_{j\in I}$ et $(u_j^-)_{j\in I}$ sont sommables.

    On définit la somme de la famille $(u_j)_{i \in I}$ par :
    $$\sum_{i \in I} u_i = \sum_{i \in I} u_i^+ - \sum_{i \in I} u_i^-$$
\end{theoreme}
\marginenote{Remarque}{Cette définition est bien compatible avec la def 5.17 dans le
cas où $\forall j\in\mathbb{B}, U_j\geqslant 0$.}
\begin{proof}{}{}

    $\Rightarrow$ On a $\left\{\begin{array}{cc}
        |u_j^+| \leqslant |u_j|\\
        |u_j^-| \leqslant |u_j|
    \end{array}\right.$, d'où l'implication par la proposition 5.16.\\

    $\Leftarrow$ Supposons que $(u_j^+)_{j\in I}$ et $(u_j^-)_{j\in I}$ sont sommables. Soit $J\subset I$ fini non vide

    $$\sum_{j\in J} |u_j|\leqslant\sum_{j\in J} U_j^+ + \sum_{j\in J} U_j^-\leqslant \sum_{j\in I} U_j^+ + \sum_{j\in I} U_j^- \leftarrow\text{constante indépendante de J}$$

    \end{proof}



\begin{theoreme}{Somme d'une famille sommable de complexes}{}

    Soit $(u_j)_{j \in I}$ une famille de complexes. Alors :
    
    $(u_j)_{j\in I}$ est sommable ssi $(\re(u_j))_{j\in I}$ et $(\im(u_j)_{j\in I}$ le sont.

    On définit la somme de la famille $(u_j)_{j \in I}$ par :
    $$\sum_{j \in I} u_j = \sum_{j \in I} \re(u_j) + i \sum_{j \in I} \im(u_j)$$

\end{theoreme}

\idee{Utiliser la proposition 5.16.}
\begin{proof}{}{}
    $\Rightarrow$
    $\begin{array}{cc}
    |\re(u_j)| \leqslant |u_j| \\
    |\im(u_j)| \leqslant |u_j|
    \end{array}$

    $\Leftarrow
    |u_j|\leqslant |\re(u_j)| + |\im(u_j)|$

\end{proof}

\begin{proposition}{Linéarité de la somme}{}
    Soit $(u_j)_{j \in I}, (v_j)_{j \in I}$ deux familles sommables de complexes et $\alpha, \beta \in \C$.

    Alors $(\alpha u_j + \beta v_j)_{j \in I}$ est sommable et :
    $$\sum_{j \in I} (\alpha u_j + \beta v_j) = \alpha \sum_{j \in I} u_j + \beta \sum_{j \in I} v_j$$
\end{proposition}

\begin{proof}
    Le faire d'abord pour des réels positifs, puis des réels quelconques, puis des complexes.
\end{proof}

\begin{theoreme}{Sommation par paquets}{}
    Soient $I$ et $K$ des ensembles non vides et $\displaystyle I = \bigsqcup_{k \in K} I_k$ union
    disjointe d'ensembles non vides.

    Soit $(u_j)_{j \in I}$ une famille de complexes. Alors $(u_j)_{j \in I}$ est sommable ssi
    
    $\forall k \in K, \, (u_j)_{j \in I_k}$ est sommable et $\displaystyle \left( \sum_{j \in I_k} |u_j| \right)_{k \in K}$ est sommable
    et en ce cas, 
    $$ \sum_{j \in I} u_j = \sum_{k \in K} \left( \sum_{j \in I_k} u_j \right)$$
\end{theoreme}

\subsection{Conséquence : permutation des termes}

\begin{theoreme}{Invariance par bijection}{}
    Soit $(u_j)_{j \in I}$ une famille sommable de complexes et $\sigma: I \to I$ une bijection 
    alors $(u_{\sigma(j)})_{j \in I}$ est sommable et $\displaystyle \sum_{j \in I} u_j = \sum_{j \in I} u_{\sigma(j)}$.
\end{theoreme}

\begin{proof}
    Posons $K = I$ et $\forall k \in K, \, I_k = \{\sigma(k)\}$.

    $\sigma$ étant une bijection, $\displaystyle I = \bigsqcup_{k \in K} I_k$.

    $(u_j)_{j \in I_k}$ étant sommable, par le théorème de sommation par paquets, 
    $\forall k \in K, \, (u_j)_{j \in I_k}$ est sommable 
    et comme $\displaystyle \sum_{j \in I_k} |u_j| = |u_{\sigma(k)}|$,
    on a que $(u_{\sigma(k)})_{k \in K}$ est sommable, et :
    $$\sum_{j \in I} u_j = \sum_{k \in K} \left( \sum_{j \in I_k} u_j \right) = \sum_{k \in K} u_{\sigma(k)} = \sum_{j \in I} u_{\sigma(j)}$$
\end{proof}

\subsection{Lien avec la dénombrabilité}

\begin{definition}{Support d'une famille}{}
    Soit $(u_j)_{j \in I}$ une famille de complexes. On appelle support de cette 
    famille l'ensemble $\{j \in I ~|~ u_j \neq 0 \}$.
\end{definition}

\begin{proposition}{Somme sur le support}{}
    Soit $(u_j)_{j \in I}$ une famille sommable de complexes de support $S$. Alors :
    $$\sum_{j \in I} u_j = \sum_{j \in S} u_j$$
\end{proposition}


\begin{proof}
    Posons $K = \{1, 2\}$, et $I_1 = S, \, I_2 = I \setminus S$.

    Par le théorème de sommation par paquets, $\displaystyle \sum_{j \in I} u_j = \sum_{j \in S} u_j + \sum_{j \in I \setminus S} u_j$.
\end{proof}

\begin{theoreme}{Nature du support d'une famille sommable}{}
    Le support d'une famille sommable est au plus dénombrable. 
\end{theoreme}

\begin{proof}
    Soit $(u_j)_{j \in I}$ une famille sommable de complexes. 

    Posons $I_0 = \{j \in I, \, 1 < |u_j| \}$, et pour $n \in \N^*$,
    $I_n = \{j \in I, \, \frac 1 {n + 1} < |u_j| \leqslant \frac 1 n \}$.

    Montrons que $\displaystyle S = \bigsqcup_{n \in \N} I_n$. 

    $\subset$ $\forall n \in \N, \, I_n \subset S$ car $\forall j \in I_n, \, u_j \neq 0$
    Si $n \neq n', \, I_n \cap I_{n'} = \emptyset$ donc $\displaystyle \bigsqcup_{n \in \N} I_n \subset S$.

    $\supset$ $j \in S$, alors $\mid u_j \mid > 0$, on a donc $\displaystyle j \in I_{\left\lfloor \frac 1 {|u_j|} \right\rfloor}$, 
    donc $S \subset \displaystyle \bigsqcup_{n \in \N} I_n$.

    Soit $k \in \N$, alors $\forall j \in I_k, \, \frac 1 {k + 1} < \underbrace{|u_j|}_{\text{sommable}}$

    Donc $\left( \frac 1 {k + 1} \right)_{j \in I_k}$ est sommable, or $\frac 1 {k + 1}$ est constant et positif, 

    donc $I_k$ est fini, donc $S$ est au plus dénombrable.
\end{proof}

\section{Familles indexées par \texorpdfstring{$\N$}{N}}
\subsection{Théorème}

\begin{theoreme}{Familles sommables indexées par $\N$}{}
    Soit $(u_n)_{n \in N}$ une famille de complexes. Alors
    cette famille est sommable ssi $\displaystyle \sum_{n \geqslant 0} u_n$ est 
    absolument convergente et alors :
    $$\sum_{n \in \N} u_n = \sum_{n = 0}^{+\infty} u_n$$    
\end{theoreme}

\begin{proof}
    $\Rightarrow$ Supposons que $(u_n)_{n \in \N}$ est sommable. 

    Ainsi $\forall J \subset \N$ fini, $\displaystyle \sum_{n \in J} |u_n| \leqslant M$.

    Soit $N \in \N$, alors $\displaystyle \sum_{j = 0 }^N |u_j|  = \sum_{n \in \{0, \ldots, N\}} |u_n| \leqslant M$.

    Donc $\sum_{j \geqslant 0} |u_j|$ converge et la série est absolument convergente. 

    \lvspace $\Leftarrow$ Supposons que $\sum_{n \geqslant 0} |u_n|$ converge. 

    Soit $J \subset \N$ fini, et soit $N \in \N$ un majorant de $J$. 

    $$\sum_{j \in J} |u_j| \leqslant \sum_{j = 0}^N |u_j| \leqslant \sum_{j = 0}^{+\infty} |u_j|$$

    Donc la famille est sommable. \\

    \uline{Si $\forall n, \, u_n \in \R^+$ :}
    Pour $J \subset \N$ fini, soit $N$ un majorant de $J$. 
    $$\sum_{j \in J} u_j \leqslant \sum_{j = 0}^N u_j \leqslant \sum_{j = 0}^{+\infty} u_j$$

    Par définition de la borne sup, 
    $$\sum_{n \in \N} u_n \leqslant \sum_{j = 0 }^{+\infty} u_j$$

    Par ailleurs $\forall k \in \N, \, \displaystyle \sum_{i = 0}^k u_i  = \sum_{i \in \ninter 0 n} u_i \leqslant \sum_{i \in \N} u_i$.

    Donc par passage à la limite, $\displaystyle \sum_{i = 0}^{+\infty} u_i \leqslant \sum_{i \in \N} u_i$ et on a donc l'égalité. \\

    \uline{Si $\forall n \in \N, \, u_n \in \R$ :}
    $a \leqslant u_n^+ \leqslant |u_n|$ et $0 \leqslant u_n^- \leqslant |u_n|$.

    Donc par ce qui précède, 
    $$\sum_{n = 0}^\infty u_n^+ = \sum_{n \in \N} u_n^+ \text{ et } \sum_{n = 0}^{+\infty} u_n^- = \sum_{n \in \N} u_n^-$$

    Et donc par défférence, $\displaystyle \sum_{n = 0}^{+\infty} u_n = \sum_{n \in \N} u_n$.

    Dans le cas complexe, on a l'égalité pour les parties réelles et imaginaires, d'où la conclusion 
    par combinaison linéaire. 
\end{proof}

\subsection{Conséquence}

\begin{proposition}{Sommation par paquets pairs et impairs - LHP}{}
    Soit $\sum_{n \geqslant 0} u_n$ une série absolument convergente, alors :
    $$\sum_{n  = 0 }^{+\infty} u_n = \sum_{p = 0 }^{+\infty} u_{2p} + \sum_{p = 0}^{+\infty} u_{2p + 1}$$
\end{proposition}

\begin{proof}
    La famille est sommable, et on applique le théorème de sommation par paquets. 
\end{proof}

\begin{exemple}
    Calcul de $\displaystyle S = \sum_{n = 0 }^{+\infty} \frac 1 {(2n + 1)^2}$, qui existe car le terme général est équivalent à $\frac 1 {n^2}$. 

    $\displaystyle \zeta(2) = \frac{\pi^2} 6 = \sum_{n = 1}^{+\infty} \frac 1 {n^2}$ famille sommable. 

    $\displaystyle \zeta(2) = \sum_{n = 0 }^{+\infty} \frac 1 {(2n)^2} + \sum_{n = 0 }^{+\infty} \frac 1 {(2n + 1)^2}$ par le th de sommation par paquets. 

    $\displaystyle \frac{\pi^2} 6 = S + \frac{\pi^2}{24}$

    Donc $S = \frac {\pi^2} 8$.
\end{exemple}

\subsection{Commutative convergence}

\begin{theoreme}{Invariance par bijection de $\N$}{}
    Soit $\displaystyle \sum_{n \geqslant 0} u_n$ une série absolument convergente et $\sigma : \N \to \N$ une bijection.
    Alors $\displaystyle \sum_{n =\geqslant 0} u_{\sigma(n)}$ est absolument convergente et :
    $$\sum_{n = 0}^{+\infty} u_n = \sum_{n = 0}^{+\infty} u_{\sigma(n)}$$
\end{theoreme}

\begin{proof}
    $(u_n)_{n \in \N}$ est sommable car la série converge absolument,
    on réutilise donc le résultat prouvé avant.
\end{proof}

\begin{remarque}
    Si la série est semi-convergente, c'est faux.
\end{remarque}


\section{Famillles indexées par \texorpdfstring{$\N^2$}{N²}}
\subsection{Cas des familles réelles positives}

\begin{theoreme}{}{}
    Soit $(u_{n, p})_{(n, p) \in \N^2}$ une famille de \textbf{réels positifs}.
    Alors les assertions suivantes sont équivalentes :
    \begin{enumerate}
        \item la famille $(u_{n, p})_{(n, p) \in \N^2}$ est sommable
        \item $\ds \forall n \in \N, \, \sum_{p \geqslant 0 } u_{n, p}$ converge et $\ds \sum_{n \geqslant 0} \left(\sum_{p = 0}^{+\infty} u_{n, p} \right)$ converge.
        \item $\ds \forall p \in \N, \, \sum_{n \geqslant 0 } u_{n, p}$ converge et $\ds \sum_{p \geqslant 0} \left(\sum_{n = 0}^{+\infty} u_{n, p} \right)$ converge.
        \item En posant $\ds \forall k \in \N, \, w_k = \sum_{n = 0 }^{k }u_{n, k - n} = \sum_{p = 0 }^{k } u_{k - p, p} = \sum_{\substack{(n, p) \in \N^2 \\ n + p = k}} u_{n, p}$, la série $\ds \sum_{k = 0}^{+\infty} w_k$ converge
            et en ce cas $\ds \sum_{(n, p) \in \N^2} u_{n,p} = \sum_{n = 0}^{+\infty} \sum_{p = 0}^{+\infty}u_{n, p} = \sum_{p = 0}^{+\infty} \sum_{n = 0}^{+\infty}u_{n, p} = \sum_{k = 0 }^{+\infty} w_k$ 
    \end{enumerate}
\end{theoreme}

\begin{proof}
    C'est le théorème de sommation par paquets appliqué aux partitions de
    $\N^2 = \bigsqcup_{n \in \N} I_n = \bigsqcup_{p \in \N} J_p = \bigsqcup_{k \in \N} E_k$
    où $I_n = \{ (n, p), \, p \in \N \}, \, J_p = \{ (n, p), \, n \in \N \}$ et
    $E_k = \{ (n, p), \, n + p = k \}$.
\end{proof}

\subsection{Théorème de Fubini positif}

\begin{theoreme}{Théorème de Fubini positif}{}
    Avec la convention qu'une somme infinie non convergente de réels positifs vaut 
    $+ \infty$ on a pour toute famille de \textbf{réels positifs} $(u_{n, p})_{(n, p) \in \N^2}$ :
    $$\sum_{n = 0}^{+\infty} \left(\sum_{p = 0}^{+\infty} u_{n, p} \right) = \sum_{p = 0}^{+\infty} \left( \sum_{n = 0}^{+\infty} u_{n, p} \right) $$
\end{theoreme}

\begin{proof}
    Conséquence du 1. de la proposition précédente.
\end{proof}

\subsection{Cas général}

\begin{theoreme}{}{}
    Soit $(u_{n, p})_{(n, p) \in \N^2}$ une famille sommable de complexes. Alors :

    \begin{enumerate}
        \item $\ds \forall n \in \N, \, \sum_{p \geqslant 0 } u_{n, p}$ converge et $\ds \sum_{n \geqslant 0} \left( \sum_{p = 0}^{+\infty} u_{n, p} \right)$ converge.
        \item $\ds \forall p \in \N, \, \sum_{n \geqslant 0 } u_{n, p}$ converge et $\ds \sum_{p \geqslant 0} \left( \sum_{n = 0}^{+\infty} u_{n, p} \right)$ converge.
        \item En posant $\ds \forall k \in \N, \, w_k = \sum_{n = 0 }^{k }u_{n, k - n} = \sum_{p = 0}^{k} u_{k - p, p} = \sum_{\substack{(n, p) \in \N^2 \\ n + p = k}} u_{n, p}$, la série $\ds \sum_{k = 0}^{+\infty} w_k$ converge
            et en ce cas $\ds \sum_{(n, p) \in \N^2} = \sum_{n = 0 }^{+\infty} \sum_{p = 0}^{+\infty}u_{n, p} = \sum_{p = 0}^{+\infty} \sum_{n = 0}^{+\infty}u_{n, p} = \sum_{k = 0}^{+\infty} w_k$ 
    \end{enumerate}
\end{theoreme}

\subsection{Produit de Cauchy}

\begin{theoreme}{Produit de Cauchy}{}
    Soient $ \sum_{n \geqslant 0} u_n$ et $ \sum_{n \geqslant 0} v_n$ deux séries 
    numériques \textbf{absolument convergentes}.
    
    On définit la suite $(w_n)_{n \in \N}$ appelée produit de convolution des suites 
    $(u_n)_{n \in \N}$ et $(v_n)_{n \in \N}$ par :
    $$\forall n \in \N, \, w_n = \sum_{\substack{(p, q) \in \N^2 \\ p + q = n}} u_p v_q = \sum_{k = 0}^n u_k v_{n - k}$$

    On définit le produit de Cauchy des séries $ \sum_{n \geqslant 0} u_n$ et $ \sum_{n \geqslant 0} v_n$ 
    comme étant la série $ \sum_{n \geqslant 0} w_n$. Alors cette série est
    absolument convergente, et :
    $$\sum_{n = 0 }^{+\infty} w_n = \left( \sum_{n = 0 }^{+\infty} u_n \right) \left( \sum_{n = 0 }^{+\infty} v_n \right)$$
\end{theoreme}

\idee{56 théorèmes de sommation par paquets.}

\begin{proof}
    Considérons le complexe 

    \begin{align*}
        P & = \sum_{n = 0 }^{+\infty} u_n \sum_{p = 0 }^{+\infty} v_p \\
            & = \sum_{p = 0 }^{+\infty} \left( \left( \sum_{n = 0}^{+\infty} u_n \right) v_p \right) \\
            & = \sum_{p = 0 }^{+\infty} \left( \sum_{n = 0}^{+\infty} (u_n v_p) \right) \\ 
    \end{align*}

    Posons $\forall n, p \in \N, \, a_{n, p} = u_n v_p$.

    Les séries étant absolument convergentes, on a l'existence de 
    $P' = \ds \sum_{n = 0 }^{+\inft } |u_n| \cdot \sum_{p = 0 }^{+\infty} |v_p|$.

    et donc de $P' = \ds \sum_{n = 0 }^{+\infty} \sum_{p = 0 }^{+\infty} |a_{n, p}|$.

    Donc, par le théorème de sommation par paquets, $(a_{n, p})_{(n, p) \in \N^2}$ est sommable. 

    Encore par le théorème de sommation par paquets, 
    $$P' = \sum_{(n, p) \in \N^2} |a_{n, p}| = \sum_{n = 0 }^{+\infty} z_n$$
    où $\ds z_n = \sum_{p = 0}^n |u_p| |v_{n - p}|$.

    Or $\forall n \in \N, \, |w_n| \leqslant z_n$, donc $\ds \sum_{n = 0 }^{+\infty} w_n < A$. 

    De plus, de nouveau par le théorème de sommation par paquets,
    $$P = \sum_{(n, p) \in \N^2} a_{n, p} = \sum_{n = 0 }^{+\infty} w_n$$
\end{proof}

\begin{exemple}
    Soient $z_1, z_2 \in \C$. Alors 
    $\displaystyle e^{z_1} = \sum_{n = 0 }^{+\infty} \frac{z_1^n}{n!}$ et
    $\displaystyle e^{z_2} = \sum_{n = 0 }^{+\infty} \frac{z_2^n}{n!}$ sont absolument convergentes.

    Et donc par produit de Cauchy, $\displaystyle e^{z_1} e^{z_2} = \sum_{n = 0 }^{+\infty} w_n$ 
    série absolument convergente où $\displaystyle w_n = \sum_{ p = 0}^n \frac{z_1^p}{p!} \frac{z_2^{n - p}}{(n - p)!}$.

    Donc $w_n = \displaystyle \sum_{p = 0 }^{n } \frac{1}{n!} \binom n p z_1^p z_2^{n - p} = \frac{(z_1 + z_2)^n}{n!}$ 
    par le binôme de Newton.

    Donc $\displaystyle e^{z_1} e^{z_2} = e^{z_1 + z_2}$.

\end{exemple}

\begin{corollaire}{Produit d'exponrntielles de matrices}{}
    Soient $A, B \in \mathscr M_p(\C)$ telles que $AB = BA$. 
    Alors :
    $$\displaystyle e^A e^B = e^{A + B}$$
\end{corollaire}

\begin{proof}
    Munissons $\mathscr M_k(\C)$ d'une norme d'algèbre $||\cdot||$. 

    Posons $\displaystyle \forall n \in \N, \, M_n = \sum_{k = 0 }^{n } \frac{A^k}{k!} \cdot \sum_{k = 0 }^{n } \frac{B^k}{k!} - \sum_{k = 0 }^{n }\frac{(A + B)^k}{k!}$.

    Alors $M_n \limit n {+ \infty} e^A \cdot e^B - e^{A + B}$ car le produit matriciel est continu
    car en dimension finie. 

    Montrons que $M_n \limit n {+ \infty} 0$ et on pourra conclure par unicité de la limite. 

    \begin{align*}
        ||M_n|| & = \left|\left| \sum_{k = 0 }^{n } \sum_{p = 0 }^{n } \frac{A^k B^p}{k! p!} - \sum_{q = 0}^n \sum_{i = 0 }^{q } \binom q i \frac{A^i B^{q - i}}{q!} \right|\right| \\
                & = \left|\left| \sum_{(k, p) \in \ninter 0 n^2} \frac{A^k B^p}{k! p!} - \sum_{q = 0}^n \sum_{k = 0}^q \frac{A^k B^{q - k}}{k! (q - k)!} \right|\right| \\
                & = \left|\left| \sum_{(k, p) \in \ninter 0 n^2} \frac{A^k B^p}{k! p!} - \sum_{\substack{(k, p) \in \ninter 0 n^2 \\ k + p \leqslant n}} \frac{A^k B^p}{k! p!} \right|\right| \\
    \end{align*}

    En effet,

    \begin{align*}
        \sum_{q = 0}^n \sum_{k = 0}^q \frac{A^k B^{q - k}}{k! (q - k)!} & = \sum_{q= 0}^n \sum_{k = 0}^n \mathbb 1_{k \leqslant q} \frac{A^k B^{q - k}}{k! (q - k)!} \\
         & = \sum_{k = 0}^n \sum_{q = 0}^n \mathbb 1_{k \leqslant q} \frac{A^k B^{q - k}}{k! (q - k)!} \\
    \end{align*}

    \begin{align*}
        \sum_{q = 0}^n \sum_{k = 0}^q \frac{A^k B^{q - k}}{k! (q - k)!} & = \sum_{ k = 0}^n \sum_{p = -k}^{n - k} \mathbb 1_{0 \leqslant p} \frac{A^k B^p}{k! p!} \\
         & = \sum_{ k = 0}^n \sum_{p = 0}^{n - k} \frac{A^k B^p}{k! p!} \\
         & = \sum_{\substack{(k, p) \in \ninter 0 n^2 \\ k + p \leqslant n}} \frac{A^k B^p}{k! p!} \\
    \end{align*}

    \begin{align*}
        ||M_n|| & = \left|\left| \sum_{\substack{(k, p) \in \ninter 0 n^2 \\ k + p > n}} \frac{A^k B^p}{k! p!} \right|\right| \\
                & \leqslant \sum_{\substack{(k, p) \in \ninter 0 n^2 \\ k + p > n}} \frac{||A||^k ||B||^p}{k! p!} \\
                & \leqslant \sum_{(k, p) \in \ninter 0 n^2} \frac{||A||^k ||B||^p}{k! p!} \text{ puis en inversant matrice et norme } \\
                & \leqslant \sum_{k = 0 }^{n} \frac{||A||^k}{k!} \cdot \sum_{p = 0 }^{n } \frac{||B||^p}{p!} - \sum_{q = 0}^n  \frac{(||A || + ||B||)^q}{q!} \\
                & \limit n {+ \infty} e^{||A||} e^{||B||} - e^{||A|| + ||B||} = 0
    \end{align*}
\end{proof}
